\documentclass{article}
\usepackage{enumitem}
\usepackage{amsmath}
\begin{document}

\title{Chapter 20: Locomotion and Movement}
\date{}
\maketitle

\begin{enumerate}[label=Q\arabic*.]

\item The sliding filament theory of muscle contraction states that:
\\(A) Myosin filaments shorten during contraction
\\(B) Actin filaments shorten during contraction
\\(C) Actin filaments slide over myosin filaments; sarcomere shortens but filament lengths remain unchanged
\\(D) Both actin and myosin filaments shorten simultaneously

\item The functional unit of skeletal muscle contraction is the:
\\(A) Myofibril
\\(B) Sarcomere (from one Z-line to the next)
\\(C) Muscle fibre
\\(D) Motor unit

\item Troponin-tropomyosin complex regulates muscle contraction by:
\\(A) Providing energy via ATP hydrolysis
\\(B) Blocking the myosin-binding sites on actin in relaxed muscle; Ca$^{2+}$ binding to troponin shifts tropomyosin to expose binding sites
\\(C) Generating the power stroke directly
\\(D) Pumping Ca$^{2+}$ into the sarcoplasmic reticulum

\item Which zone of the sarcomere contains only myosin (thick filaments)?
\\(A) I band
\\(B) H zone (within the A band)
\\(C) Z line
\\(D) M line only

\item Assertion: ATP is required both for the power stroke and for the relaxation of muscle.\\
Reason: ATP is hydrolyzed to cock the myosin head and is also needed by the Ca$^{2+}$-ATPase pump to return Ca$^{2+}$ to the SR.
\\(A) Both true, Reason is correct explanation
\\(B) Both true, Reason is not correct explanation
\\(C) Assertion true, Reason false
\\(D) Assertion false, Reason true

\item Rigor mortis (stiffening of muscles after death) occurs because:
\\(A) Actin filaments are destroyed
\\(B) ATP depletion prevents detachment of myosin heads from actin
\\(C) Troponin irreversibly binds Ca$^{2+}$
\\(D) Sarcoplasmic reticulum releases all Ca$^{2+}$ permanently

\item Which type of bone joint allows the greatest range of movement?
\\(A) Fibrous joint (suture)
\\(B) Cartilaginous joint (symphysis)
\\(C) Ball-and-socket synovial joint
\\(D) Hinge joint

\item The human skeleton has how many bones in an adult?
\\(A) 300
\\(B) 270
\\(C) 206
\\(D) 180

\item Which of the following is the correct match of bone type and example?
\\(A) Long bone: Carpals
\\(B) Short bone: Femur
\\(C) Flat bone: Scapula
\\(D) Irregular bone: Humerus

\item The neuromuscular junction transmits impulse from nerve to muscle via:
\\(A) Norepinephrine
\\(B) Acetylcholine (ACh) binding to nicotinic receptors on motor end plate
\\(C) Direct electrical coupling
\\(D) GABA

\item Osteoporosis is caused by:
\\(A) Excess calcium deposition in bone matrix
\\(B) Decreased bone mineral density due to imbalance between osteoblast and osteoclast activity
\\(C) Vitamin D excess causing bone thickening
\\(D) Cartilage degeneration in joints

\item The I band of a sarcomere:
\\(A) Contains only thick (myosin) filaments
\\(B) Contains only thin (actin) filaments; bisected by Z line
\\(C) Contains both thick and thin filaments
\\(D) Disappears completely during contraction

\item Which mineral is released from the sarcoplasmic reticulum to trigger muscle contraction?
\\(A) Na$^+$
\\(B) K$^+$
\\(C) Ca$^{2+}$
\\(D) Mg$^{2+}$

\item Match the following muscle disorders with their descriptions:\\
1. Myasthenia gravis $\rightarrow$ (a) Autoimmune destruction of ACh receptors at NMJ\\
2. Muscular dystrophy $\rightarrow$ (b) Genetic deficiency of dystrophin protein\\
3. Tetanus $\rightarrow$ (c) Sustained muscle contraction due to \textit{Clostridium tetani} toxin\\
4. Gout $\rightarrow$ (d) Uric acid crystal deposition in joints
\\(A) 1-a, 2-b, 3-c, 4-d
\\(B) 1-b, 2-a, 3-d, 4-c
\\(C) 1-c, 2-d, 3-a, 4-b
\\(D) 1-d, 2-c, 3-b, 4-a

\item The axial skeleton includes:
\\(A) Pectoral girdle, pelvic girdle, and limb bones
\\(B) Skull, vertebral column, sternum, and ribs
\\(C) Femur, humerus, and their associated girdles
\\(D) Only the skull and vertebral column

\end{enumerate}
\end{document}
